\documentclass[12pt]{article}
\usepackage{amsmath,amssymb,amsthm}
\usepackage[margin=1in]{geometry}
\usepackage{hyperref}

\newtheorem{theorem}{Theorem}
\newtheorem{lemma}[theorem]{Lemma}
\newtheorem{corollary}[theorem]{Corollary}
\newtheorem{remark}{Remark}

\newcommand{\R}{\mathbb{R}}
\newcommand{\E}{\mathbb{E}}
\newcommand{\Zmat}{\mathcal{Z}}
\newcommand{\bx}{\mathbf{x}}
\newcommand{\by}{\mathbf{y}}
\newcommand{\bv}{\mathbf{v}}
\newcommand{\bd}{\mathbf{d}}

\title{Correctness of the DE-MCz Slice Sampler\\with Frozen Archive and Capped Iterations}
\author{}
\date{}

\begin{document}
\maketitle

We prove that the DE-MCz Slice Sampler with a frozen reference archive
and fixed-iteration (capped) slice sampling preserves the target
distribution~$\pi$. The proof proceeds in five steps: (1)~the direction
law is symmetric and independent of the current state, (2)~standard
slice sampling preserves~$\pi$, (3)~the capped variant also
preserves~$\pi$, (4)~the single-walker kernel preserves~$\pi$, and
(5)~the fully parallel sweep preserves~$\pi^N$.

Throughout, $\pi$ denotes a target density on $\R^d$ with respect to
Lebesgue measure, and $\Zmat = \{\bx^{(1)}, \ldots, \bx^{(M)}\}
\subset \R^d$ is a fixed (frozen) reference archive of size~$M \geq 2$.

%% ====================================================================
\section*{Theorem 1: Direction Symmetry and Independence}

\begin{theorem}[Direction law]
\label{thm:direction}
Define the direction law $q_\Zmat$ as follows: draw indices $i, j$
uniformly from $\{1, \ldots, M\}$ with $i \neq j$, and set
\[
    \bv = \frac{\bx^{(i)} - \bx^{(j)}}{\|\bx^{(i)} - \bx^{(j)}\|}.
\]
Then:
\begin{enumerate}
    \item[(a)] $q_\Zmat(\bv) = q_\Zmat(-\bv)$  \quad (symmetry).
    \item[(b)] The direction $\bv$ is independent of the current walker
    state~$\bx$, conditional on $\Zmat$.
\end{enumerate}
\end{theorem}

\begin{proof}
(a)~The pair $(i, j)$ is drawn uniformly from all $M(M-1)$ ordered
pairs with $i \neq j$. Swapping $i \leftrightarrow j$ sends $\bv
\mapsto -\bv$ and is a bijection on the set of ordered pairs. Hence
$q_\Zmat(\bv) = q_\Zmat(-\bv)$.

(b)~The draw of $(i, j)$ depends only on the PRNG key and $M$, not on
the current state~$\bx$. Since $\Zmat$ is frozen, the direction is
determined entirely by $\Zmat$ and the PRNG key.
\end{proof}

%% ====================================================================
\section*{Theorem 2: Standard Slice Sampling Invariance}

\begin{theorem}[Neal, 2003]
\label{thm:slice}
Let $\pi(\bx)$ be a target density on $\R^d$, and let $\bd \in \R^d$
be a fixed unit direction. The one-dimensional slice sampling procedure
along~$\bd$ --- consisting of (i)~drawing a slice height
$u \sim \mathrm{Uniform}(0, \pi(\bx))$, (ii)~finding the slice
$S = \{t \in \R : \pi(\bx + t\bd) > u\}$, and (iii)~drawing
$t^* \sim \mathrm{Uniform}(S)$ to obtain $\bx' = \bx + t^*\bd$ ---
leaves $\pi$ invariant.
\end{theorem}

\begin{proof}
This is the standard result from Neal~(2003, Theorem~1). The augmented
target $p(\bx, u) = \mathbf{1}[0 < u < \pi(\bx)]$ is jointly uniform
on the region under $\pi$. The Gibbs update of $\bx$ given $u$ draws
uniformly from the level set $\{\bx : \pi(\bx) > u\}$ restricted to
the line through $\bx$ in direction $\bd$, which preserves the
augmented distribution. Marginalizing over $u$ recovers $\pi$.
\end{proof}

\begin{remark}
In practice, the slice $S$ is found approximately via stepping-out
(expanding the bracket $[L, R]$ until $\pi(\bx + L\bd) \leq u$ and
$\pi(\bx + R\bd) \leq u$) followed by shrinking (proposing uniformly
in $[L, R]$ and narrowing the bracket on rejection). Neal~(2003) shows
that the stepping-out/shrinking procedure also preserves~$\pi$, by a
careful pathwise pairing argument: each accepted trajectory $\bx
\to \bx'$ is paired with a reverse trajectory $\bx' \to \bx$ that
encounters the same sequence of bracket adjustments.
\end{remark}

%% ====================================================================
\section*{Theorem 3: Capped Slice Sampling Invariance}

This is the main novel result. Our implementation caps the shrinking
phase at $N_{\mathrm{shrink}}$ iterations. If no acceptable point is
found within the budget, the chain performs a self-transition
($\bx' = \bx$).

\begin{theorem}[Capped slice kernel]
\label{thm:capped}
Let $K_{\mathrm{cap}}(\bx, \cdot)$ denote the Markov kernel obtained
by running the slice sampling shrinking procedure for at most
$N_{\mathrm{shrink}}$ iterations. If an acceptable point is found
within budget, return it; otherwise return $\bx' = \bx$. Then
$K_{\mathrm{cap}}$ leaves $\pi$ invariant:
\[
    \int \pi(\bx)\, K_{\mathrm{cap}}(\bx, A)\, d\bx = \pi(A)
    \quad \text{for all measurable } A.
\]
\end{theorem}

\begin{proof}
We decompose the capped kernel as a state-dependent mixture of two
$\pi$-invariant kernels.

For a given state $\bx$, slice height $u$, and bracket $[L, R]$, let
$p_{\mathrm{find}}(\bx, u, L, R)$ denote the probability that the
shrinking procedure finds an acceptable point within
$N_{\mathrm{shrink}}$ iterations. Conditional on finding a point, the
accepted proposal is drawn from the slice $S \cap [L, R]$ via the
standard shrinking procedure, which preserves $\pi$ by Neal's pathwise
argument (Theorem~\ref{thm:slice} and the Remark above).

Let $K_{\mathrm{exact}}$ denote the exact (uncapped) shrinking kernel
and $K_{\mathrm{id}}(\bx, A) = \mathbf{1}[\bx \in A]$ the identity
kernel. Both preserve $\pi$: $K_{\mathrm{exact}}$ by Neal's proof, and
$K_{\mathrm{id}}$ trivially.

The capped kernel is:
\[
    K_{\mathrm{cap}}(\bx, A)
    = p_{\mathrm{find}}(\bx) \cdot K_{\mathrm{exact}}(\bx, A)
    + (1 - p_{\mathrm{find}}(\bx)) \cdot K_{\mathrm{id}}(\bx, A).
\]
The key observation: the pathwise argument still holds for trajectories
that terminate within budget. For any accepted trajectory of length
$k \leq N_{\mathrm{shrink}}$, the reverse trajectory (from the
accepted point back to the original) has the \emph{same} length~$k$,
because the shrinking bracket adjustments are symmetric about the
origin $t = 0$ (the original point's projection). Therefore:
\begin{itemize}
    \item If the forward trajectory accepts in $k \leq N_{\mathrm{shrink}}$
    steps, the reverse also accepts in $k$ steps. Both transitions are
    handled by $K_{\mathrm{exact}}$.
    \item If the forward trajectory does \emph{not} accept in
    $N_{\mathrm{shrink}}$ steps, the reverse trajectory from $\bx$
    (which is returned unchanged) would also not accept in
    $N_{\mathrm{shrink}}$ steps, since the trajectory is $\bx \to \bx$
    (the identity). Both transitions are handled by $K_{\mathrm{id}}$.
\end{itemize}

Since $\pi K_{\mathrm{exact}} = \pi$ and $\pi K_{\mathrm{id}} = \pi$,
we have for any measurable $A$:
\begin{align*}
    \int \pi(\bx)\, K_{\mathrm{cap}}(\bx, A)\, d\bx
    &= \int \pi(\bx)\, p_{\mathrm{find}}(\bx)\, K_{\mathrm{exact}}(\bx, A)\, d\bx \\
    &\quad + \int \pi(\bx)\, (1 - p_{\mathrm{find}}(\bx))\, K_{\mathrm{id}}(\bx, A)\, d\bx.
\end{align*}
This is \emph{not} simply the sum of two $\pi$-invariant kernels
weighted by a state-dependent function (which would not generally
preserve $\pi$). Instead, the decomposition works because the
pathwise pairing ensures that $K_{\mathrm{exact}}$ in the capped
regime --- i.e., restricted to trajectories of length
$\leq N_{\mathrm{shrink}}$ --- still preserves $\pi$ when paired
with the identity on the complementary event. The pairing ensures
the detailed balance relation holds trajectory-by-trajectory.
\end{proof}

\begin{remark}
The stepping-out phase is also capped at $N_{\mathrm{expand}}$
iterations per side. This does not affect invariance: a narrower
bracket $[L, R] \subseteq [L', R']$ simply restricts the set of
attainable proposals. The shrinking procedure within any bracket
containing the origin ($t = 0$, corresponding to $\bx' = \bx$)
preserves $\pi$, because the identity transition is always available as
a fallback.
\end{remark}

%% ====================================================================
\section*{Theorem 4: Single-Walker Kernel}

\begin{theorem}[Single-walker invariance]
\label{thm:walker}
Let $K_\Zmat(\bx, \cdot)$ denote the single-walker kernel: draw
direction $\bv \sim q_\Zmat$, then apply the capped slice sampling
kernel $K_{\mathrm{cap}}$ along $\bv$. Then $K_\Zmat$ leaves $\pi$
invariant.
\end{theorem}

\begin{proof}
For each fixed direction $\bv$, the kernel
$K_{\mathrm{cap}}^{(\bv)}(\bx, \cdot)$ preserves $\pi$
(Theorem~\ref{thm:capped}). The single-walker kernel marginalizes
over directions:
\[
    K_\Zmat(\bx, A) = \int K_{\mathrm{cap}}^{(\bv)}(\bx, A)\,
    q_\Zmat(\bv)\, d\bv.
\]
Since $\pi K_{\mathrm{cap}}^{(\bv)} = \pi$ for every $\bv$,
integrating over $\bv$ preserves this:
\[
    \int \pi(\bx)\, K_\Zmat(\bx, A)\, d\bx
    = \int \left[\int \pi(\bx)\, K_{\mathrm{cap}}^{(\bv)}(\bx, A)\,
    d\bx\right] q_\Zmat(\bv)\, d\bv
    = \int \pi(A)\, q_\Zmat(\bv)\, d\bv = \pi(A).
\]
\end{proof}

%% ====================================================================
\section*{Theorem 5: Parallel Sweep}

\begin{theorem}[Parallel sweep invariance]
\label{thm:sweep}
Let $(\bx_1, \ldots, \bx_N)$ be the positions of $N$ walkers. Define
the sweep kernel $K_\Zmat^{\mathrm{sweep}}$ as the product of $N$
independent applications of $K_\Zmat$, all conditioned on the same
frozen archive $\Zmat$:
\[
    K_\Zmat^{\mathrm{sweep}}(\bx_{1:N}, A_{1:N})
    = \prod_{i=1}^N K_\Zmat(\bx_i, A_i).
\]
Then $K_\Zmat^{\mathrm{sweep}}$ leaves the product target
$\pi^N(\bx_{1:N}) = \prod_{i=1}^N \pi(\bx_i)$ stationary.
\end{theorem}

\begin{proof}
With $\Zmat$ frozen, walker $i$'s update depends only on $\bx_i$, the
PRNG key, and $\Zmat$. It does \emph{not} depend on any other walker's
state $\bx_j$ ($j \neq i$). This is the critical distinction from
standard ensemble methods (e.g., Goodman \& Weare 2010), which require
complementary-set splitting because directions are drawn from other
walkers' current positions.

Since each $K_\Zmat$ preserves $\pi$ (Theorem~\ref{thm:walker}), and
the walkers are conditionally independent given $\Zmat$, the product
kernel preserves $\pi^N$:
\[
    \int \pi^N(\bx_{1:N})\, K_\Zmat^{\mathrm{sweep}}(\bx_{1:N},
    A_{1:N})\, d\bx_{1:N}
    = \prod_{i=1}^N \int \pi(\bx_i)\, K_\Zmat(\bx_i, A_i)\, d\bx_i
    = \prod_{i=1}^N \pi(A_i) = \pi^N(A_{1:N}).
\]
\end{proof}

%% ====================================================================
\section*{Discussion}

\paragraph{Adaptive warmup.}
During the warmup phase, the archive $\Zmat$ grows as new walker
positions are appended. This makes the algorithm adaptive in the sense
of Roberts \& Rosenthal (2007): the transition kernel changes between
sweeps. The production phase freezes $\Zmat$, eliminating this
adaptive component and yielding a time-homogeneous Markov chain for
which the above theorems apply directly. This is the simplest route to
rigorous guarantees; proving ergodicity of the adaptive warmup phase
would require verifying diminishing adaptation and containment
conditions, which we avoid by freezing.

\paragraph{Mu tuning.}
The slice width parameter $\mu$ is tuned during warmup based on the
observed bracket expansion ratio. Since $\mu$ is fixed before the
production phase begins, the production kernel is fully specified and
the invariance proofs apply without modification.

\paragraph{Ergodicity.}
The theorems above establish that $\pi$ (resp.\ $\pi^N$) is a
stationary distribution of the production kernel. Ergodicity ---
convergence to $\pi$ from any starting distribution --- requires
additionally that the kernel is irreducible and aperiodic. Slice
sampling is aperiodic (the identity transition has positive probability
whenever $N_{\mathrm{shrink}}$ is finite). Irreducibility holds
provided the archive $\Zmat$ spans $\R^d$ in the sense that the
direction vectors $\{(\bx^{(i)} - \bx^{(j)}) / \|\bx^{(i)} -
\bx^{(j)}\| : i \neq j\}$ span all of $\R^d$, which is satisfied
whenever $\Zmat$ contains $d + 1$ points in general position.

%% ====================================================================
\paragraph{References.}
\begin{itemize}
    \item Neal, R.~M. (2003). Slice Sampling. \textit{Annals of Statistics}, 31(3), 705--767.
    \item ter Braak, C.~J.~F. \& Vrugt, J.~A. (2008). Differential Evolution Markov Chain
          with snooker updater and fewer chains. \textit{Statistics and Computing}, 18(4), 435--446.
    \item Karamanis, M. \& Beutler, F. (2021). Ensemble Slice Sampling: Parallel, black-box
          and gradient-free inference for correlated \& multimodal distributions.
          \textit{arXiv:2002.06212}.
    \item Goodman, J. \& Weare, J. (2010). Ensemble samplers with affine invariance.
          \textit{Comm.\ Appl.\ Math.\ Comp.\ Sci.}, 5(1), 65--80.
    \item Roberts, G.~O. \& Rosenthal, J.~S. (2007). Coupling and ergodicity of adaptive
          Markov chain Monte Carlo algorithms. \textit{J.\ Appl.\ Probab.}, 44(2), 458--475.
\end{itemize}

\end{document}
